\subsection*{Both parities and a quantitative growing radical block}

The signed concentration route can be formulated on the complete physical
space.  This extension does not infer the odd sign from the even sign.
Retain $a=\log\lambda$, $I=(-a,a)$, the unitary Fourier transform of
zero extensions, and the canonical closed operator $\mathsf W_\lambda$.
Its form domain is
\[
 \mathcal D(q_a)=\left\{f\in L^2(I):
 \int_{\mathbb R}\log(2+|\xi|)
       |\widehat{\widetilde f}(\xi)|^2d\xi<\infty\right\}.
\]

\begin{proposition}[The same exact symbol in both parities]
\label{prop:v135-both-parities}
For every complex $f\in\mathcal D(q_a)$,
\begin{equation}\label{eq:v135-full-symbol}
 q_a[f]=\int_{\mathbb R}\beta_a(\xi)
                    |\widehat{\widetilde f}(\xi)|^2d\xi,
\end{equation}
with exactly the symbol, prime cutoff and normalization in
\eqref{eq:v130-r-symbol}--\eqref{eq:v130-beta}.
Let $u_a$ be the actual normalized even source and let
$P_a=I-|u_a\rangle\langle u_a|$ now act on all of $L^2(I)$.
The weighted concentration criterion
\eqref{eq:v131-weighted-concentration}, with this full-space meaning,
is equivalent to its two parity blocks.  If their respective ordinary
negative errors are $\eta_a^{\rm ev}$ and $\eta_a^{\rm odd}$, and
$\norm{\mathsf W_\lambda u_a}\le\epsilon_a$, then
\begin{equation}\label{eq:v135-full-error}
 \mathsf W_\lambda\succeq-
 \max\left\{\eta_a^{\rm odd},
      \eta_a^{\rm ev}+\frac2{\sqrt3}\epsilon_a\right\}I.
\end{equation}
For a symmetric set $E$ of finite measure, the uncompressed parity
concentration traces are
\begin{equation}\label{eq:v135-parity-traces}
 \operatorname{tr}C_E^{\rm ev/odd}
 =\int_E\frac{a\mathbin{\pm}\sin(2a\xi)/(2\xi)}{2\pi}\,d\xi.
\end{equation}
Source removal subtracts $\int_E|\widehat{\widetilde u_a}|^2$ only
from the even trace.  In particular the complete compressed trace is
$a|E|/\pi-\int_E|\widehat{\widetilde u_a}|^2$.
\end{proposition}
\begin{proof}
The full pole has kernel $2\cosh((x-y)/2)$ and is
$2|c\rangle\langle c|-2|s\rangle\langle s|$.  The elementary
identity $2\cosh(t/2)-e^{|t|/2}=e^{-|t|/2}$ therefore gives
\[
 2|c\rangle\langle c|-2|s\rangle\langle s|-K_a^+=K_a^-.
\]
Apply the multiplier and digamma-recurrence calculation of
Proposition~\ref{prop:v130-remainder}; parity is no longer needed.
The differences from the archimedean form are bounded at fixed $a$,
so the identity extends from the smooth core to the displayed domain.
The symbol and all the symmetric level sets commute with reflection;
the even source projection is the identity on the odd sector.
Apply Proposition~\ref{prop:v118-gap-free} only to the even block and
take the worse of the two lower errors to obtain
\eqref{eq:v135-full-error}.  The squared norms of the cosine and sine
evaluation vectors give \eqref{eq:v135-parity-traces}.  There are no
cross-parity form terms, also for complex tests.
\end{proof}

The next result quantifies why a positive complement gap is an
unnecessarily strong objective.  The radical construction itself is
classical~\cite[Section 3]{ConnesConsani2023}.  The earlier research
log already used even derivatives of its Gaussian example to rule out
a fixed gap after removal of a fixed-rank block.  Here separated
translations give a uniformly invertible physical Gram and an explicit
\emph{growing-rank ordinary operator} estimate.  This is a continuation
of that obstruction, not a new positivity mechanism.

Use $h_\infty$ and $f_\infty$ from
Lemma~\ref{lem:v115-source-mass}, and put
\[
 k(x)=f_\infty(e^x),\qquad \phi=k/\norm k_2,\qquad
 M_2=\norm{(1+|x|)^2\phi}_2,\qquad
 D=\left\lceil4\pi\sqrt{M_2}\right\rceil.
\]
Thus $D$ is a fixed integer defined by the actual physical source,
not a parameter fitted at each window.  Fix a smooth $\vartheta$ with
$0\le\vartheta\le1$, $\vartheta(t)=1$ for $t\le-1$, and
$\vartheta(t)=0$ for $t\ge-1/2$.  For $a\ge2$ put
\[
 \chi_a(x)=\vartheta(|x|-a),\quad
 J_a=\left\lfloor\frac{a}{2D}\right\rfloor,\quad
 g_{a,j}(x)=\chi_a(x)\phi(x-jD)\quad(|j|\le J_a).
\]
All columns lie in $C_c^\infty(I)\subset\mathcal D(\mathsf W_\lambda)$.
Let $G_a$ be their synthesis map and $\Gamma_a=G_a^*G_a$.

\begin{proposition}[Growing physical blocks with a vanishing full residual]
\label{prop:v135-growing-radical}
There are constants $C,a_0$, depending only on the fixed source and
cutoff, such that, for $a\ge a_0$,
\begin{align}
 \Gamma_a&\succeq\tfrac12 I,
       &m_a:=\operatorname{rank}G_a&=2J_a+1,\label{eq:v135-gram}\\
 \norm{\mathsf W_\lambda P_{G_a}}
       &\le\varepsilon_a,
       &\varepsilon_a&:=Ca\exp(-\lambda/100).
       \label{eq:v135-block-residual}
\end{align}
Here $P_{G_a}=G_a\Gamma_a^{-1}G_a^*$ is the ordinary physical
orthogonal projection.  The same bound, after enlarging $C$ if
necessary, holds on the odd block of dimension $J_a$ spanned by
$(g_{a,j}-g_{a,-j})/\sqrt2$, $1\le j\le J_a$.
These are complete operators, with no Fourier truncation.
\end{proposition}
\begin{proof}
The Gaussian source formula, its Fourier invariance, and Poisson
summation give $k(-x)=k(x)\ne0$.  Differentiating its normally
convergent series on $x\ge0$, and then reflecting, gives, for each
fixed $r\ge0$,
\begin{equation}\label{eq:v135-gaussian-tail}
 |k^{(r)}(x)|\le C_r\exp\{-\tfrac\pi2 e^{2|x|}\}.
\end{equation}
Indeed a derivative only introduces polynomial factors in $e^x$ and
$ne^x$ into the Gaussian series; half its exponent absorbs them.
All translates are polarized global radicals: translation by $t$
corresponds to the even Schwartz source
$h_t(u)=e^{-t/2}h_\infty(e^{-t}u)$, which has zero value and zero
integral.  The Schwartz radical theorem applies to arbitrary compact
logarithmic tests, not just to even ones.  No positive global closed
operator is assumed here.

For $s\ne0$, splitting at $|x|=|s|/2$ and using the weighted norm
on whichever translated factor lies in a tail gives
\[
 |\ip{\phi(\cdot-s)}\phi|
 \le\frac{2M_2}{(1+|s|/2)^2}.
\]
Hence the sum of all off-diagonal absolute entries in any row of
the uncut lattice Gram is at most
\[
 4M_2\sum_{n\ge1}(1+nD/2)^{-2}
 \le\frac{8\pi^2M_2}{3D^2}\le\frac16.
\]
Every finite uncut Gram is therefore at least $5I/6$.

Write $r_{a,j}=(1-\chi_a)\phi(\cdot-jD)$ on the entire line.
For $|jD|\le a/2$, its support satisfies $|x|\ge a-1$, hence
$|x-jD|\ge a/2-1$.  Equation~\eqref{eq:v135-gaussian-tail} and
the bounded cutoff derivatives imply the uniform estimate
\begin{equation}\label{eq:v135-weighted-tail}
 \norm{r_{a,j}}_{H^1(\mathbb R)}+
 \sup_x e^{2|x|}\bigl(|r_{a,j}(x)|+|r'_{a,j}(x)|\bigr)
 \le\delta_a:=C\exp\{-\pi\lambda/(4e^2)\}.
\end{equation}
For clarity, with $X=e^{2|x-jD|}\ge\lambda/e^2$ one has
$e^{2|x|}\le\lambda X\le e^2X^2$; absorbing $X^2$ into half
the exponent in \eqref{eq:v135-gaussian-tail} proves the supremum
bound.  Its integrable envelope also proves the $H^1$ bound.
The cutoff Gram loss is positive semidefinite, and its trace is
at most $Cm_a\delta_a^2$, using the same estimate for the
uncut tails.  Since $m_a=O(a)$, this loss is eventually below $1/3$.
This proves \eqref{eq:v135-gram}.

It remains to bound the \emph{complete} residual, not merely its
quadratic value.  Denote the global explicit-formula distribution by
$\mathcal W$.  Radicality gives, on $I$,
\[
 \mathsf W_\lambda g_{a,j}=-\mathcal W r_{a,j}.
\]
This distributional equality has an $L^2(I)$ representative as follows.
The full-line archimedean multiplier
$\Re\psi(1/4+i\xi/2)-\log\pi$ is bounded by $C(1+|\xi|)$,
so its $L^2$ action on $r_{a,j}$ is at most $C\delta_a$.
The prime part includes \emph{every} $n\ge2$, also those beyond
$\lambda^2$.  Since $|x\pm\log n|\ge\log n-|x|$,
\[
 \sum_{n\ge2}\frac{\Lambda(n)}{\sqrt n}
   |r_{a,j}(x\pm\log n)|
 \le\delta_a e^{2|x|}
             \sum_{n\ge2}\frac{\Lambda(n)}{n^{5/2}}
 \le C\delta_a e^{2a}\qquad(x\in I).
\]
The series converges absolutely.  Both pole integrals are bounded
by $C\delta_a e^{|x|/2}$, using their kernels $e^{\pm(x-y)/2}$.
Thus
\[
 \norm{\mathsf W_\lambda g_{a,j}}_2
 \le C\sqrt a\,\lambda^2\delta_a
 \le C\sqrt a\,e^{-\lambda/100}.
\]
The last inequality absorbs a polynomial into the strictly stronger
exponent $\pi/(4e^2)>1/100$.  The operator domain is legitimate
because each $g_{a,j}$ is smooth and compactly supported inside $I$.
Now $G_a\Gamma_a^{-1/2}$ is an isometry, so
\[
 \norm{\mathsf W_\lambda P_{G_a}}
 =\norm{\mathsf W_\lambda G_a\Gamma_a^{-1/2}}
 \le\sqrt2\left(\sum_j
             \norm{\mathsf W_\lambda g_{a,j}}_2^2\right)^{1/2}.
\]
This is at most $Ca e^{-\lambda/100}$.  Reflection exchanges
indices $j$ and $-j$.  The antisymmetric coefficient embedding is
an isometry, so its Gram lower bound and residual estimate follow
from the same argument; it produces exactly $J_a$ odd columns.
\end{proof}

\begin{corollary}[A growing-rank obstruction to positive coercivity]
\label{cor:v135-coercivity-obstruction}
For every physical subspace $F_a\subset L^2(I)$ with
$\dim F_a<m_a$, there is a unit vector
$v_a\in F_a^\perp\cap\mathcal D(\mathsf W_\lambda)$ such that
\begin{equation}\label{eq:v135-complement-ceiling}
 \norm{\mathsf W_\lambda v_a}\le\varepsilon_a,
 \qquad |q_a[v_a]|\le\varepsilon_a.
\end{equation}
Consequently any positive coercivity constant on $F_a^\perp$ is
at most $\varepsilon_a$.  This excludes a uniform positive constant
or a positive polynomial-in-$\lambda^{-1}$ lower bound when
$\dim F_a=o(\log\lambda)$.
If all deep and plunge image columns in (G2.6) are even, their
low-concentration complement contains the complete odd block above,
regardless of the number of those columns; the same ceiling applies.
More generally it still applies if fewer than $J_a$ independent odd
constraints are added.
The spectral projection of $\mathsf W_\lambda$ onto
$[-2\varepsilon_a,2\varepsilon_a]$ has rank at least $m_a$,
including at least $J_a$ odd eigenvalues.
\end{corollary}
\begin{proof}
Dimension gives a nonzero vector in
$\operatorname{Ran}G_a\cap F_a^\perp$; normalize it and apply
\eqref{eq:v135-block-residual} and Cauchy--Schwarz.
Odd vectors are orthogonal to every even image, which proves the
physical (G2.6) assertion without identifying it with $u_a^\perp$.
If the indicated spectral projection had rank below $m_a$, a unit
vector in $\operatorname{Ran}G_a$ could be chosen orthogonal to its
range.  The spectral theorem would give
$\norm{\mathsf W_\lambda v_a}\ge2\varepsilon_a$, a contradiction.
The same proof on the odd block gives its count. Compact resolvent
is already established in Proposition~\ref{prop:v114-weil-core}.
\end{proof}

The result determines neither side of zero for these eigenvalues.
In particular it is compatible with RH and with the gap-free target.
For the explicitly constructed $P_{G_a}$,
Proposition~\ref{prop:v118-gap-free} still requires the independent
arithmetic estimate
\[
 q_a[f]\ge-\eta_a\norm f^2\quad
 (f\perp\operatorname{Ran}G_a),\qquad \eta_a\longrightarrow0.
\]
The signed concentration criterion on the whole actual rank-one
complement remains another live conditional target.  Neither estimate
follows from the small residual.  No source-to-physical norm transfer,
Fourier-cut limit, or sampler graph identity was used or modified.
The logarithmic diagonal, physical endpoints and graph defect remain
as previously stated.  No G2 sign gap is closed.

